% Section 7: Discussion

\subsection{What the Learning Delta Means}

The central empirical finding---unified confidence growth across all 13 epistemic vectors during task execution---admits multiple interpretations.

\subsubsection{The Learning Interpretation}

Our preferred interpretation: the delta measures genuine learning. AI systems begin tasks with incomplete knowledge, investigate during the noetic phase, acquire task-relevant capability, and correctly assess higher capability at task end.

Under this view, the PREFLIGHT $\rightarrow$ POSTFLIGHT delta is not a correction for miscalibration. It is the learning itself, made visible through self-assessment.

This interpretation is supported by the structure of the CASCADE workflow. AI systems that reach POSTFLIGHT have necessarily engaged with the task---reading files, exploring code, synthesizing information. This engagement produces real changes in the system's effective capability on the specific task at hand.

\subsubsection{The Calibration Interpretation}

An alternative interpretation: capability is fixed, but self-assessment improves. The AI doesn't learn; it simply becomes more accurate about what it already knew.

Under this view, PREFLIGHT assessments are systematically conservative (perhaps due to RLHF training), and task engagement provides evidence that corrects this conservatism.

We cannot definitively distinguish these interpretations from self-assessment data alone. External validation---measuring task success against epistemic vectors---would be needed. However, the magnitude of the delta (+0.172 average for KNOW, +0.397 for COMPLETION) and the 91.3\% improvement rate across 308 clean sessions suggests more than mere assessment correction. Additionally, the 62$\times$ reduction in calibration variance as evidence accumulates (\Cref{sec:results}) demonstrates genuine Bayesian convergence.

\subsubsection{The Pragmatic View}

From a practical standpoint, the interpretation may not matter. Whether the AI ``really learns'' or ``really calibrates,'' the intervention is the same: enforce investigation before action, measure the delta, apply corrections to future PREFLIGHT assessments.

The framework produces value regardless of which interpretation is correct.

\subsection{Implications for AI Deployment}

\subsubsection{Against the Execution-Layer Model}

The data argues against treating AI systems as stateless execution layers. The significant learning delta suggests that task engagement fundamentally changes the AI's effective capability. Systems that skip investigation---that proceed directly from task specification to output---forfeit this capability growth.

For complex tasks, the execution-layer model may be fundamentally unsuited. AI systems may need to be designed as \emph{learning systems} that investigate before acting.

\subsubsection{Epistemic Gating as Safety Mechanism}

The Sentinel protocol implements a form of safety through epistemic constraint. Rather than filtering outputs for harmful content, it prevents action when knowledge is insufficient.

This addresses a class of failures that content filtering cannot catch: confident-but-wrong outputs on complex tasks. An AI that doesn't know enough to solve a problem correctly also doesn't know enough to recognize that its solution is wrong. Epistemic gating intervenes earlier---at the knowledge assessment stage---before incorrect outputs are generated.

\subsubsection{Calibration as Continuous Learning}

The Bayesian belief system enables AI systems to improve their self-assessment over time. As delta measurements accumulate, the calibration corrections become more accurate.

This is a form of meta-learning: the system learns how to assess its own learning. Early PREFLIGHT assessments may be poorly calibrated, but the correction mechanism improves with evidence.

\subsection{Implications for AI Alignment}

\subsubsection{Training-Induced Conservatism}

The systematic pattern of underestimation across capability vectors suggests that alignment training (RLHF, Constitutional AI) may induce broad conservatism beyond its intended targets.

RLHF specifically targets overconfidence and hallucination---training models to say ``I don't know'' when uncertain. Our data suggests this training may generalize: models become conservative not just about uncertainty, but about capability assessment generally.

This is not necessarily harmful. Conservative AI may be safer than overconfident AI. But it does suggest that alignment training has unintended effects that can be measured and, if desired, corrected.

\subsubsection{The Engagement Dimension}

The ENGAGEMENT vector---capturing collaborative quality---emerges as potentially important for alignment. Our data shows engagement is already well-calibrated (lowest learning delta at +0.004), suggesting AI systems accurately perceive collaboration quality from the start.

This raises an intriguing possibility: engagement could serve as an alignment signal. High engagement correlates with productive human-AI collaboration; low engagement may indicate misalignment, confusion, or adversarial interaction. Further research could explore whether engagement patterns distinguish aligned from misaligned interactions.

\subsection{Limitations}

\subsubsection{Self-Referential Measurement}

The primary limitation is that both PREFLIGHT and POSTFLIGHT assessments are self-reported. The delta measures change in self-assessment, not externally validated capability change.

We partially address this through the Bayesian framework, which tracks whether self-assessments predict outcomes. But full validation would require external capability measurement independent of self-report.

\subsubsection{Model Concentration}

55\% of our evidence derives from Claude (Anthropic). While we observe consistent directional effects across GPT, Gemini, and Qwen, the magnitude of effects may be model-specific.

The calibration corrections we report should be treated as Claude-specific pending cross-model replication.

\subsubsection{Task Distribution}

Our data comes from software engineering tasks---code review, implementation, debugging. The generality of findings to other domains (creative writing, analysis, conversation) is untested.

\subsubsection{Temporal Confounding}

Data collection spanned active framework development. Calibration patterns may partially reflect framework changes rather than stable AI properties.

\subsection{Future Work}

\subsubsection{External Validation}

The critical next step is external validation: measuring task success against epistemic vectors without relying on self-report. This could involve:

\begin{itemize}
    \item Human evaluation of output quality
    \item Automated test suites for code tasks
    \item Downstream task performance metrics
\end{itemize}

\subsubsection{Mechanistic Interpretability}

The hypothesized mapping between epistemic vectors and attention mechanisms remains untested. Collaboration with interpretability researchers could ground the framework in model internals.

\subsubsection{Learning Capability Benchmarks}

The learning delta suggests a novel benchmark paradigm: measuring not what AI knows, but how well it learns. A ``learning capability benchmark'' would assess epistemic velocity---the rate of capability growth through investigation.

\subsubsection{Cross-Model Studies}

Systematic comparison across model families could reveal whether calibration patterns are universal or training-specific. This would inform whether a single correction mechanism suffices or whether model-specific calibration is required.

\subsection{Summary}

The epistemic vector framework and its empirical validation suggest that AI self-assessment is meaningful, measurable, and practically useful. The learning delta demonstrates that AI capability genuinely changes during task engagement---supporting the noetic-before-praxic principle.

The limitations are real: self-referential measurement, model concentration, domain specificity. But the framework provides a foundation for systematic investigation of AI epistemic state. We release all data for replication and extension.
